\chapter{Introduction} \label{sec:introduction}
\section{Problem Declaration} \label{ch:problem_declaration}
Modern AI systems that work with user images must do two things reliably: keep the raw images safe, organized, quickly retrievable; run repeatable processing on those images, e.g., annotation and embedding (see Section \ref{sec:image_embeddings}\footnote{Terms like embeddings are explained later in the image processing chapter of the thesis (\ref{sec:image_processing}).}). In practice, this typically requires building a data storage that persists a large number of small binary files. In addition to storage, data processing pipelines are constructed to execute jobs on data and process them as needed.

The first operational aspect of image processing systems is the design of the storage system. The problem of storing a large number of small binary files (images) does not fit the requirements for which traditional storage solutions were designed. Storing a large number of small files can lead to significant overhead in metadata management, which can cause poor performance on traditional systems.\cite{beaver_finding_2010}

For an image store, it is crucial to minimize metadata management overhead, enable parallel read (and possibly write) operations to serve data efficiently, and focus on systems that prioritize reading over writing and deletion. At the same time, customer and regulatory constraints often rule out the public cloud, requiring full control and on-premises deployment. Among other needs, the solution must typically be fault-tolerant and highly available.

On the computation side, there are mixed requirements from online and offline processing. For offline processing, it is necessary to ensure fault tolerance and limit the number of errors during computation. Offline processing often computes large amounts of data in one bulk operation, in case of backfills or recomputation when a model changes. Bulk operations require careful scheduling of resources to deliver the largest throughput possible, but keep the price of such computation reasonably low. On the other hand, online computations focus on low latency and high availability. This leads to an auto-scaling computation solution within the hard limits of an on-premises cluster. The processing pipeline must be idempotent and resumable in the event of failure, enabling failure recovery. Finally, computed artifacts need to be persistently stored and indexed with metadata and a backlink to its original source.

\section{Motivation} \label{sec:motivation}
The motivation for this thesis is driven by a significant and growing industry need to build reliable image-processing pipelines that can operate at a massive scale. Modern machine learning applications, particularly in domains such as e-commerce, recommendation systems, and search, increasingly depend on the ability to efficiently process and analyze large image datasets. Providing a solution that can ingest, store and process millions of items is a challenge to serving the needs of large, data-intensive systems.

This thesis is directly motivated by a real-world business requirement from Recombee, a company specializing in AI-powered recommendation systems. Such systems should enable the company to develop improved recommendations and AI-powered search systems. Such a system is a key component in supporting more sophisticated content-based recommendations and search functionalities that rely on a deep understanding of visual data, rather than just provided metadata.

The challenge of efficiently managing and storing a massive volume of small binary files, often referred to as \textit{the Small File Problem}, is not unique to the primary use case of this thesis. It is a common challenge in many other data-intensive industries. For example, large-scale e-commerce platforms must manage billions of product images, and social media networks must store vast amounts of user-generated photos. Therefore, the results and research on the storage component of the thesis can be valuable for other use cases when developing such on-premises storage infrastructure, particularly for small binary files.

The second part of the thesis (developing a data processing pipeline) can also serve as inspiration for other implementations of similar use cases, where a data processing pipeline is required. Such systems could be similar to AI-processing, model training, etc.

\section{Purpose} \label{sec:purpose}
The purpose of this thesis is to design and implement an on-premises infrastructure that makes large-scale image processing possible. In practice, the work aims to provide a fault-tolerant storage solution for handling large numbers of small images ($\approx$10 KiB to 1 MiB) and to build an automated, on-demand job-execution pipeline that scales ML operations, such as embedding computations. This involves first analyzing the specific storage and processing requirements for such datasets, then researching and deploying a suitable scalable, fault-tolerant, distributed open-source storage infrastructure, and finally developing and deploying the necessary pipeline.